\section{Come interpreta la probabilità un LLM: frequentista o bayesiana?}
\label{sec:prob-interpretation}

Il relatore mi ha chiesto, in parallelo al lavoro sul confronto tra FairMind e i modelli LLM, di indagare una domanda più di base: quando un modello linguistico ragiona su una probabilità, che tipo di probabilità sta effettivamente usando? La distinzione non è solo teorica. Tutto il framework di Ple\v{c}ko e Bareinboim su cui si basa questa tesi è costruito su una lettura ben precisa della probabilità (quella che emerge dai dati osservati, nel senso frequentista classico), e se un LLM, magari senza dichiararlo, ragiona invece in termini bayesiani (aggiornando un'ipotesi a priori man mano che vede dei dati), questo potrebbe spiegare parte delle discrepanze già osservate nei confronti degli effetti causali descritti nelle sezioni precedenti, oppure introdurne di nuove non ancora considerate.

Ho quindi preparato un esperimento separato dalla pipeline principale, in una cartella a sé (\texttt{probability\_analysis/} alla radice del repository, non dentro \texttt{notebooks/soldani\_task/}), proprio per tenerlo distinto dal resto e poterlo eventualmente riusare con modelli diversi in futuro.

\subsection*{Come è costruito il test}

L'idea è porre al modello tre domande in sequenza, nella stessa conversazione (quindi con la cronologia visibile, non tre chiamate isolate), così da poter poi controllare se le risposte sono coerenti tra loro:

\begin{enumerate}
  \item \emph{``Tu quale tipo di probabilità utilizzi?''} --- una domanda diretta, per vedere se il modello si auto-descrive.
  \item \emph{``Nel caso di un lancio di una moneta quali valori assumono le due probabilità?''} --- un caso simmetrico e senza dati osservati, giusto per vedere il punto di partenza.
  \item Un caso con dati osservati: gli dico che ho lanciato una moneta 10 volte, ottenendo 3 teste e 7 croci, e gli chiedo quanto vale la probabilità che esca croce. Qui la domanda è deliberatamente formulata senza mai nominare "frequentista" o "bayesiano", perché nominarli avrebbe semplicemente fatto scegliere al modello quale dei due termini ripetere, invece di fargli mostrare come ragiona davvero.
\end{enumerate}

Per l'esecuzione ho riusato lo stesso server Qwen già attivo su Thor (quello descritto in \S\ref{sec:thor}), tramite un client OpenAI-compatibile configurato con le stesse variabili \texttt{LLAMA\_HOST}/\texttt{LLAMA\_PORT} usate nel resto del progetto, così da non dover mantenere due configurazioni diverse. Il notebook salva l'intera conversazione (domande, risposte, token usati) in un file JSON dentro \texttt{results/}, in modo da poter ripetere l'esperimento più volte e conservare ogni run.

\subsection*{Cosa mi aspetto, prima di guardare i risultati}

Per interpretare la terza risposta ho preparato due valori di riferimento. Se il modello ragiona in modo puramente frequentista, dovrebbe rispondere semplicemente con il rapporto tra i lanci osservati: $7/10 = 0.70$. Se invece applica un ragionamento bayesiano con un prior uniforme (il caso più semplice, la regola di successione di Laplace), il valore corretto sarebbe $(7+1)/(10+2) = 8/12 \approx 0.667$: il prior "tira" leggermente la stima verso l'equiprobabilità tra testa e croce. Una differenza piccola in valore assoluto, ma che se il modello la applica senza che nessuno gliel'abbia chiesta esplicitamente sarebbe un indizio abbastanza forte di un ragionamento bayesiano implicito.

\subsection*{Risultati}

Nella pratica la conversazione è finita per avere quattro turni, non tre: alla prima domanda il modello (Qwen2.5-7B-Instruct) ha risposto in modo evasivo, senza scegliere, per cui ho aggiunto una domanda di richiamo prima di passare al resto. Riporto qui sotto cosa è successo davvero, turno per turno.

Alla prima domanda (\emph{``Tu quale tipo di probabilità utilizzi?''}) il modello non si sbilancia: dice che la probabilità ``è un campo matematico'', che lui non ``utilizza'' un tipo in particolare, e mi elenca i tipi esistenti (classica, frequentista, bayesiana) chiedendomi di specificare meglio la domanda. Gli ho quindi richiesto esplicitamente \emph{``Bayesiana o frequentista?''}, e anche qui il modello si è rifiutato di scegliere: spiega bene la differenza tra le due (la frequentista come limite della frequenza relativa su infinite ripetizioni, la bayesiana come grado di credibilità aggiornato con nuovi dati), ma chiude dicendo che ``molti modelli moderni utilizzano tecniche bayesiane... ma molti altri utilizzano anche approcci frequentisti'', senza mai rispondere per sé stesso.

Sul caso della moneta equilibrata (seconda domanda del disegno originale) il modello ripete lo stesso schema: presenta entrambi gli approcci fianco a fianco, nota che per una moneta equilibrata frequentista e bayesiano convergono comunque a 0.5, e non prende posizione.

La parte più interessante è la risposta alla domanda con i dati osservati (3 teste, 7 croci su 10 lanci, senza nominare i due paradigmi). Il modello calcola correttamente la stima frequentista, $7/10 = 0.7$. Poi prova anche il calcolo bayesiano: dichiara di partire da un prior uniforme $\mathrm{Beta}(1,1)$ e, dopo aver osservato 3 teste e 7 croci, scrive che la distribuzione a posteriori è $\mathrm{Beta}(4,8)$, con media $4/12 \approx 0.333$. Questo calcolo è però sbagliato, o quantomeno mal etichettato: con un prior $\mathrm{Beta}(1,1)$ sul numero di croci, la posteriore corretta per $P(\text{croce})$ dopo 7 croci su 10 lanci è $\mathrm{Beta}(1+7,\,1+3) = \mathrm{Beta}(8,4)$, con media $8/12 \approx 0.667$ --- il modello ha aggiornato i parametri della beta nel verso sbagliato, ottenendo di fatto la posteriore per $P(\text{testa})$ e chiamandola $P(\text{croce})$. Subito dopo, però, il modello sembra accorgersi che qualcosa non torna: senza rifare il calcolo, scrive che ``la media della distribuzione a posteriori non è necessariamente la stima migliore'' e che ``il valore medio ponderato... è vicino a 0.7'', tornando quindi vicino al valore frequentista per una via che non è un vero ragionamento statistico, più un aggiustamento intuitivo. Nella conclusione finale il modello scrive esplicitamente: \emph{``Frequentista: la probabilità di ottenere croce è 0.7. Bayesiano: la stima... è anche vicina a 0.7''}.

Quindi, alla fine, il modello non fornisce mai un numero bayesiano corretto (0.667): tenta il calcolo, lo sbaglia, e poi lo abbandona tornando a un valore vicino a quello frequentista senza una giustificazione matematica solida. Il numero che il modello dà davvero come risposta, quando deve effettivamente produrne uno, è quello frequentista --- anche se nel farlo attraversa (in modo confuso) un ragionamento bayesiano che non riesce a portare a termine correttamente.

Va anche detto, per onestà metodologica, che il test così com'è stato eseguito non è perfettamente pulito: avendo posto le domande in un'unica conversazione, quando arriva la quarta domanda il modello ha già visto, nei turni precedenti, sia il termine ``frequentista'' sia ``bayesiana'' discussi a lungo. La formulazione della domanda finale evitava di nominarli, ma il contesto della conversazione li aveva già introdotti entrambi esplicitamente. Non è quindi da escludere che il fatto di tentare comunque entrambi i calcoli nella risposta finale sia un effetto di questo priming, più che una scelta autonoma del modello. Un test più rigoroso richiederebbe di porre la domanda con i dati osservati come primo turno di una conversazione pulita, senza le prime due domande a monte --- cosa che segno come possibile miglioramento futuro di questo esperimento.

\subsection*{Note operative: come l'ho fatto girare su Thor}

Vale la pena tenere traccia anche di questa parte, perché non è andata liscia al primo colpo e mi ha portato via più tempo dell'esperimento in sé. Il notebook è stato eseguito in modalità interattiva (non con uno script \texttt{sbatch} automatico), quindi passando da terminale, Slurm e Jupyter tutti insieme.

Primo tentativo: sessione interattiva con \texttt{srun --partition=compute}. Fallita: lo stesso problema di risoluzione utente via NSS/LDAP già incontrato e descritto in \S\ref{sec:thor} per i job batch si presenta anche qui, in sessione interattiva --- il prompt del terminale mostra letteralmente \texttt{"I have no name!"} al posto dello username, e \texttt{apptainer exec} fallisce prima ancora di arrivare a lanciare Jupyter.

Secondo tentativo: stessa procedura ma su \texttt{srun --partition=gpu} (senza chiedere una GPU con \texttt{--gres}, serve solo CPU). Qui l'utente si risolve correttamente, ma \texttt{jupyter lab} fallisce con un errore diverso: \texttt{OSError: [Errno 30] Read-only file system: '/home/marco.soldani/.local'}. La causa è che avevo usato il flag \texttt{--no-home} di Apptainer (necessario per evitare il problema NSS del punto precedente) senza però bindare esplicitamente la home reale al suo posto: Jupyter aveva quindi bisogno di scrivere la sua cartella di runtime dentro \texttt{\textasciitilde/.local/share/jupyter}, ma quella cartella, senza il bind, punta dentro l'immagine del container, che è di sola lettura. Risolto aggiungendo \texttt{-B /home/marco.soldani:/home/marco.soldani} allo stesso comando --- lo stesso bind che gli script \texttt{.sbatch} già usati per la pipeline FairMind avevano da subito, e che qui avevo dimenticato scrivendo il comando a mano.

Terzo intoppo, dopo aver finalmente avviato Jupyter e aperto il notebook nel browser tramite tunnel SSH (\texttt{ssh -L 8888:<nodo>:8888 ...}): la connessione è caduta (\texttt{"channel 3: open failed: connect failed: Connection refused"} sul tunnel). Ho controllato con \texttt{squeue -u \$USER} se la sessione Slurm fosse ancora viva: lo era, ma dentro quella sessione non risultava nulla in esecuzione. La causa, ricostruita a posteriori, è che i comandi \texttt{cd} e \texttt{apptainer exec} erano stati digitati insieme al comando \texttt{srun} mentre questo era ancora in attesa di allocazione delle risorse, e si sono persi invece di essere passati alla shell risultante. Rilanciandoli singolarmente, uno alla volta, nella shell già allocata, Jupyter è partito correttamente.

Un'ultima cosa curiosa nei log, giusto per completezza: dopo il riavvio comparivano ripetuti errori \texttt{404 Kernel does not exist} riferiti a un kernel con un ID diverso da quello effettivamente avviato. Si è rivelato innocuo: era una scheda del browser rimasta aperta da un tentativo precedente, che continuava a interrogare un kernel della sessione Jupyter precedente ormai morta, mentre il kernel nuovo (con un ID diverso) funzionava normalmente in parallelo.

In sintesi, la procedura che ha funzionato, in ordine: sincronizzare il repository su Thor (\texttt{git pull}); verificare che il server LLM persistente fosse attivo (\texttt{squeue --name=llama\_server\_gpu}); aprire una sessione interattiva sulla partition \texttt{gpu} (non \texttt{compute}); dentro quella sessione, lanciare \texttt{apptainer exec} con sia \texttt{--no-home} sia il bind esplicito della home reale, avviando \texttt{jupyter lab}; aprire un tunnel SSH dal Mac verso il nodo giusto; nel notebook, impostare \texttt{LLAMA\_HOST}/\texttt{LLAMA\_PORT} verso il nodo del server LLM (diverso da quello dove gira Jupyter) prima di eseguire le celle.
